% 第5章 磁有序与磁结构（Order and magnetic structures）
\chapter{磁有序与磁结构}

上一章介绍了固体中磁矩之间起作用的各种磁相互作用。本章讨论这些相互作用可以产生
哪些不同类型的磁基态。其中一些基态示于图 5.1：包括所有磁矩平行排列的铁磁体，
相邻磁矩反平行排列的反铁磁体，磁矩方向随位置沿圆锥或圆周进动的螺形与螺旋结构，
以及磁矩冻结在随机取向上的自旋玻璃。本章的任务是粗线条地说明：上一章讨论的
相互作用如何导致这些不同的基态。下一章将在更一般的语境中考察有序现象，并看到
有序是对称性破缺的后果。

\begin{figure}
\centering
\includegraphics[width=0.75\textwidth]{fig5_01.pdf}
\caption{有序体系中各种自旋排列：（a）铁磁体，（b）反铁磁体，（c）自旋玻璃，
（d）螺形与（e）螺旋结构。}
\end{figure}
\mnote{记号提醒：本书中 $\mathsf{J}$ 表示交换常数，$J$ 表示总角动量。}

\section{铁磁性}

铁磁体即使在没有外场时也有自发磁化强度，所有磁矩沿唯一的同一方向排列。\fn{1}{事实上，许多铁磁样品由于磁畴的存在，这一点在样品内并不处处成立。每个畴内有均匀的磁化强度，但相邻畴的磁化方向不同。磁畴详见第 6.7 节。}这一效应一般来源于上一章描述的交换相互作用。处于外场
$\mathbf{B}$ 中的铁磁体，需求解的哈密顿量为
\begin{equation*}
\hat{\mathcal{H}} = -\sum_{ij}\mathsf{J}_{ij}\,\mathbf{S}_i\cdot\mathbf{S}_j + g\mu_{\mathrm{B}}\sum_j\mathbf{S}_j\cdot\mathbf{B},\tag{5.1}
\end{equation*}
此时最近邻交换常数为正，以保证铁磁排列。右边第一项是海森伯交换能（见式 4.7），
第二项是塞曼能量（见式 1.35）。为简单起见，先设\fn{2}{这一假设将在 5.1.4 节放宽。}所处理的是没有轨道角动量的系统，即 L=0、J=S。

\subsection{铁磁体的外斯模型}

为推进式 5.1 的求解，需要作一个近似。定义第 $i$ 个格点上的有效分子场（molecular
field）
\begin{equation*}
\mathbf{B}_{\mathrm{mf}} = -\frac{2}{g\mu_{\mathrm{B}}}\sum_j\mathsf{J}_{ij}\,\mathbf{S}_j.\tag{5.2}
\end{equation*}
\bio{Pierre Weiss\\（1865--1946）}

聚焦第 $i$ 个自旋。它的能量来自塞曼部分 $g\mu_{\mathrm{B}}\mathbf{S}_i\cdot\mathbf{B}$
与交换部分。第 $i$ 个自旋与邻居的总交换相互作用是
$-2\sum_j\mathsf{J}_{ij}\mathbf{S}_i\cdot\mathbf{S}_j$，因子 2 来自重复计数。\fn{3}{见第 4.2.1 节。}这一项可写作
\begin{equation*}
-2\mathbf{S}_i\cdot\sum_j\mathsf{J}_{ij}\mathbf{S}_j = -g\mu_{\mathrm{B}}\mathbf{S}_i\cdot\mathbf{B}_{\mathrm{mf}}.\tag{5.3}
\end{equation*}
于是交换相互作用被邻居自旋产生的有效分子场 $B_{\mathrm{mf}}$ 取代。有效哈密顿量
现在可写作
\begin{equation*}
\hat{\mathcal{H}} = g\mu_{\mathrm{B}}\sum_i\mathbf{S}_i\cdot(\mathbf{B}+\mathbf{B}_{\mathrm{mf}})\tag{5.4}
\end{equation*}
看上去就像处于磁场 $B+B_{\mathrm{mf}}$ 中的顺磁体的哈密顿量。这一做法背后的假设是：
所有磁性离子感受同样的分子场。这相当可疑——尤其在接近磁相变的温度——下一章
再谈。对铁磁体，分子场的作用是使邻居磁矩排齐，因为主导的交换相互作用为正
（反铁磁体中则为负）。

既然分子场量度系统有序化的效果，可以设
\begin{equation*}
\mathbf{B}_{\mathrm{mf}} = \lambda\mathbf{M}\tag{5.5}
\end{equation*}
$\lambda$ 是把分子场强度表示为磁化强度函数的常数。铁磁体 $\lambda>0$。由于交换
相互作用涉及巨大的库仑能，铁磁体中的分子场往往极为巨大。

现在可以把这个问题当作置于磁场 $B+B_{\mathrm{mf}}$ 中的简单顺磁体来处理。低温下，
即使完全不加外场，内部分子场也能把磁矩排齐。注意：磁矩的排齐产生了内部分子场，
而正是这个场导致排齐——这是一个“先有鸡还是先有蛋”的局面。低温下磁有序是
自维持的。温度升高，热涨落逐渐破坏磁化强度，在某个临界温度有序被摧毁。这一模型
就是铁磁性的外斯模型（Weiss model）。

求该模型的解需要联立求解
\begin{equation*}
\frac{M}{M_{\mathrm{s}}} = \mathrm{B}_J(y)\tag{5.6}
\end{equation*}
（见式 2.38）与
\begin{equation*}
y = \frac{g_J\mu_{\mathrm{B}}J(B+\lambda M)}{k_{\mathrm{B}}T}\tag{5.7}
\end{equation*}
（见式 2.37）。\mnote{提醒：此阶段我们假定 J=S、L=0。}若没有来自分子场的 $\lambda M$ 项，
这与 2.4.3 节对顺磁体的处理完全一样。

\begin{figure}
\centering
\includegraphics[width=0.45\textwidth]{fig5_02.pdf}
\caption{B=0 时式 5.6 与 5.7 的图解法。}
\end{figure}

这两个方程可以图解。先限于 B=0，此时 $M=k_{\mathrm{B}}Ty/g_J\mu_{\mathrm{B}}J\lambda$：
把 M 对 y 作图得到的直线斜率正比于温度 T（图 5.2）。高温时，除了原点（y=0、
$M_{\mathrm{s}}=0$）外式 5.6 与 5.7 无联立解。当直线斜率小于布里渊函数在原点的
斜率时，情形改变：低温下出现三个解——一个 $M_{\mathrm{s}}=0$，另两个 $M_{\mathrm{s}}$
为 ± 某非零值。可以证明：当直线不如原点处的布里渊函数陡时，非零解稳定而零解
不稳定（若 $T<T_{\mathrm{C}}$ 时系统处于 $M_{\mathrm{s}}=0$，任何涨落——无论多小
——都会使系统转入两个稳定态之一）。于是低于某一温度时出现非零磁化强度，随材料
冷却而增长：物质即使在无外场时也被磁化。这一自发磁化强度正是铁磁性的特征。

转变温度可由“直线 $M=k_{\mathrm{B}}Ty/g_J\mu_{\mathrm{B}}J\lambda M_{\mathrm{s}}$ 与曲线
$M=M_{\mathrm{s}}\mathrm{B}_J(y)$ 在原点斜率相等”的条件得到。$y$ 小时
$\mathrm{B}_J(y)=(J+1)y/3J+O(y^{3})$。转变温度即居里温度 $T_{\mathrm{C}}$：
\begin{equation*}
T_{\mathrm{C}} = \frac{g_J\mu_{\mathrm{B}}(J+1)\lambda M_{\mathrm{s}}}{3k_{\mathrm{B}}} = \frac{n\lambda\mu_{\mathrm{eff}}^{2}}{3k_{\mathrm{B}}}.\tag{5.8}
\end{equation*}
于是分子场 $B_{\mathrm{mf}}=\lambda M_{\mathrm{s}}=3k_{\mathrm{B}}T_{\mathrm{C}}/g_J\mu_{\mathrm{B}}(J+1)$；
对 $J=\frac{1}{2}$、$T_{\mathrm{C}}\sim10^{3}$ K 的铁磁体，
$B_{\mathrm{mf}}=k_{\mathrm{B}}T_{\mathrm{C}}/\mu_{\mathrm{B}}\sim1500$ T——巨大的
有效磁场，反映了交换相互作用的强度。

这些方程的解作为温度的函数绘于图 5.3（对一系列 J）。曲线形状随 J 略有不同，但
一般特征不变：$T\geq T_{\mathrm{C}}$ 时磁化强度为零，$T<T_{\mathrm{C}}$ 时非零；
在 $T=T_{\mathrm{C}}$ 处磁化强度连续但其梯度不连续。在此分子场模型中，非磁相与
铁磁相之间的相变因此属于二级相变（second-order phase transition）。相变的级定义为
自由能最低阶导数在转变处出现不连续的阶数：一级相变在自由能一阶导数（体积、熵、
磁化强度等量）处不连续跳变，熵的跳变给出潜热；二级相变在二阶导数（压缩率、热容
等量）处不连续。此处不连续出现在磁化强度的梯度，即自由能的二阶导数——故转变为
二级。相变与临界指数将在 6.4 节详述。一些常见铁磁体的性质列于表 5.1。

\begin{figure}
\centering
\includegraphics[width=0.45\textwidth]{fig5_03.pdf}
\caption{对不同 J 值导出的平均场磁化强度随温度的函数。}
\end{figure}

\begin{table}
\centering
\caption{一些常见铁磁体的性质。}
\begin{small}
\begin{tabular}{lcc}
\toprule
材料 & $T_{\mathrm{C}}$（K） & 磁矩（$\mu_{\mathrm{B}}$/化学式单元）\\
\midrule
Fe & 1043 & 2.22\\
Co & 1394 & 1.715\\
Ni & 631 & 0.605\\
Gd & 289 & 7.5\\
MnSb & 587 & 3.5\\
EuO & 70 & 6.9\\
EuS & 16.5 & 6.9\\
\bottomrule
\end{tabular}
\end{small}
\end{table}

\subsection{磁化率}

在 $T\geq T_{\mathrm{C}}$ 下加小场 B 给出小磁化强度，可用布里渊函数的 $y\ll1$ 近似：
\begin{equation*}
\frac{M}{M_{\mathrm{s}}} \approx \frac{g_J\mu_{\mathrm{B}}(J+1)}{3k_{\mathrm{B}}}\left(\frac{B+\lambda M}{T}\right)\tag{5.9}
\end{equation*}
于是
\begin{equation*}
\frac{M}{M_{\mathrm{s}}} \approx \frac{T_{\mathrm{C}}}{\lambda M_{\mathrm{s}}}\left(\frac{B+\lambda M}{T}\right).\tag{5.10}
\end{equation*}
整理得
\begin{equation*}
\frac{M}{M_{\mathrm{s}}}\left(1-\frac{T_{\mathrm{C}}}{T}\right) \approx \frac{T_{\mathrm{C}}B}{\lambda M_{\mathrm{s}}}\tag{5.11}
\end{equation*}
从而
\begin{equation*}
\chi = \lim_{B\to0}\frac{\mu_0M}{B} \propto \frac{1}{T-T_{\mathrm{C}}}\tag{5.12}
\end{equation*}
即居里--外斯定律（Curie--Weiss law）。

\begin{figure}
\centering
\includegraphics[width=0.45\textwidth]{fig5_04.pdf}
\caption{$B\neq0$ 时式 5.6 与 5.7 的图解法。}
\end{figure}

\subsection{磁场的效果}

加磁场的效果是把图解中的直线向右平移（见图 5.4）：一切温度下都有 $M\neq0$ 的解，
相变被消除。非零磁场中的铁磁体总有能量优势使磁化强度非零、磁矩沿场排列。相变的
消除见图 5.5——图中对一系列磁场绘出式 5.6 与 5.7 的图解。此模型无须考虑磁场施加
的方向：无论沿哪个方向加场，磁化强度都会转向跟随；模型中不存在铁磁体自身特有的
方向。真实铁磁体则不然——必须考虑材料的磁各向异性（见下一章）。

\begin{figure}
\centering
\includegraphics[width=0.45\textwidth]{fig5_05.pdf}
\caption{$J=\frac{1}{2}$ 时对不同的外加场 B 计算的平均场磁化强度随温度的函数。
只有 B=0 时才有相变。}
\end{figure}

在 $T=T_{\mathrm{C}}$ 处，磁场的效应可以解析求出：小场下 $M\propto B^{1/3}$。证明
需要取 $\mathrm{B}_J(y)$ 泰勒展开的下一项。写 $\mathrm{B}_J(y)=(J+1)y/3J-\zeta y^{3}+O(y^{5})$
（$\zeta$ 为常数），联立 $M=M_{\mathrm{s}}\mathrm{B}_J(y)$ 与
\begin{equation*}
y = \frac{g_J\mu_{\mathrm{B}}J(B+\lambda M)}{k_{\mathrm{B}}T_{\mathrm{C}}} = \frac{(B+\lambda M)}{\zeta M_{\mathrm{s}}}\tag{5.13}
\end{equation*}
得
\begin{equation*}
M = M_{\mathrm{s}}\zeta\frac{(B+\lambda M)}{\lambda} - \zeta M_{\mathrm{s}}\left(\frac{3J(B+\lambda M)}{\lambda(J+1)M_{\mathrm{s}}}\right)^{3}\tag{5.14}
\end{equation*}
从而
\begin{equation*}
B \propto (B+\lambda M)^{3}\tag{5.15}
\end{equation*}
由于 $\lambda M\gg B$，右侧由 $M^{3}$ 项主导，故 $M\propto B^{1/3}$。

\subsection{分子场的起源}

1907 年外斯提出分子场模型时，令他失望的是：为符合自然界中很大的 $T_{\mathrm{C}}$，
常数 $\lambda$ 必须取得很大。只考虑偶极场无法解释需要 $\sim10^{3}$ T 的内场
（如前所述，这是解释 Fe 居里温度所必需的）。三十年后，海森伯证明：正是涉及巨大
库仑能的交换相互作用，才是巨大分子场的来源。\fn{4}{分子场是一个方便的虚构——不要以为铁磁体中的电子真的感受 $\sim10^{3}$ T 的磁场。交换相互作用纯属静电效应。谈分子场时，我们是假装交换相互作用其实是内部磁场。要点是：如果它真的存在而海森伯交换不存在，那么分子场必须 $\sim10^{3}$ T 才能解释我们看到的现象。}

由 $\lambda$ 参数化的分子场可与由 $\mathsf{J}_{ij}$ 参数化的交换相互作用大小联系
起来。设交换相互作用只在一个离子的 z 个最近邻内起作用、取值 $\mathsf{J}$，则用
式 5.1、5.4、5.5 容易得到
\begin{equation*}
\lambda = \frac{2z\mathsf{J}}{ng^{2}\mu_{\mathrm{B}}^{2}}.\tag{5.16}
\end{equation*}
再用式 5.8，居里温度可写为
\begin{equation*}
T_{\mathrm{C}} = \frac{2z\mathsf{J}J(J+1)}{3k_{\mathrm{B}}}.\tag{5.17}
\end{equation*}

迄今我们假定 L=0、J=S——这对多数 3d 离子成立。交换发生在自旋自由度之间，因此
依赖 S；而离子的磁矩依赖总（自旋+轨道）角动量 J。对 3d 离子二者是同一回事，因为
L 被淬灭。

对 4f 离子则不然：S 不是好量子数，J 才是。于是 S 垂直于 J 的分量必平均为零，
平行于 J 的分量才守恒——必须把 S 投影到 J 上。由 $J=L+S$，$L+2S$ 等于 $g_JJ$
加一个垂直于 J 的分量，故 S 中作为好量子数的分量是 $(g_J-1)J$。各 4f 离子的
$(g_J-1)$ 列于表 5.2。可见对所谓的重稀土（Gd 至 Yb），S 与 J 平行；对所谓的轻稀土
（Ce 至 Sm），S 与 J 反平行。

\begin{table}
\centering
\caption{用洪特规则得到的 4f 离子 $g$ 因子。}
\begin{small}
\begin{tabular}{lccccccc}
\toprule
离子 & 壳层 & $S$ & $L$ & $J$ & $g_J$ & $g_J-1$ & $(g_J-1)^{2}J(J+1)$\\
\midrule
Ce$^{3+}$ & $4f^{1}$ & $\frac{1}{2}$ & 3 & $\frac{5}{2}$ & $\frac{6}{7}$ & $-\frac{1}{7}$ & 0.18\\
Pr$^{3+}$ & $4f^{2}$ & 1 & 5 & 4 & $\frac{4}{5}$ & $-\frac{1}{5}$ & 0.80\\
Nd$^{3+}$ & $4f^{3}$ & $\frac{3}{2}$ & 6 & $\frac{9}{2}$ & $\frac{72}{99}$ & $-\frac{27}{99}$ & 1.84\\
Pm$^{3+}$ & $4f^{4}$ & 2 & 6 & 4 & $\frac{3}{5}$ & $-\frac{2}{5}$ & 3.20\\
Sm$^{3+}$ & $4f^{5}$ & $\frac{5}{2}$ & 5 & $\frac{5}{2}$ & $\frac{2}{7}$ & $-\frac{5}{7}$ & 4.46\\
Eu$^{3+}$ & $4f^{6}$ & 3 & 3 & 0 & -- & -- & --\\
Gd$^{3+}$ & $4f^{7}$ & $\frac{7}{2}$ & 0 & $\frac{7}{2}$ & 2 & 1 & 15.75\\
Tb$^{3+}$ & $4f^{8}$ & 3 & 3 & 6 & $\frac{3}{2}$ & $\frac{1}{2}$ & 10.50\\
Dy$^{3+}$ & $4f^{9}$ & $\frac{5}{2}$ & 5 & $\frac{15}{2}$ & $\frac{4}{3}$ & $\frac{1}{3}$ & 7.08\\
Ho$^{3+}$ & $4f^{10}$ & 2 & 6 & 8 & $\frac{5}{4}$ & $\frac{1}{4}$ & 4.50\\
Er$^{3+}$ & $4f^{11}$ & $\frac{3}{2}$ & 6 & $\frac{15}{2}$ & $\frac{6}{5}$ & $\frac{1}{5}$ & 2.55\\
Tm$^{3+}$ & $4f^{12}$ & 1 & 5 & 6 & $\frac{7}{6}$ & $\frac{1}{6}$ & 1.17\\
Yb$^{3+}$ & $4f^{13}$ & $\frac{1}{2}$ & 3 & $\frac{7}{2}$ & $\frac{8}{7}$ & $\frac{1}{7}$ & 0.32\\
Lu$^{3+}$ & $4f^{14}$ & 0 & 0 & 0 & -- & -- & --\\
\bottomrule
\end{tabular}
\end{small}
\end{table}

用 $(g_J-1)\mathbf{J}$ 替换 $\mathbf{S}$ 中守恒的部分，表达式
$-\sum_{ij}\mathsf{J}_{ij}\mathbf{S}_i\cdot\mathbf{S}_j$ 可以换成
$-\sum_{ij}(g_J-1)^{2}\mathsf{J}_{ij}\mathbf{J}_i\cdot\mathbf{J}_j$。再用磁矩
$\boldsymbol{\mu}=-g_J\mu_{\mathrm{B}}\mathbf{J}$ 重复导出式 5.16 的计算，得
\begin{equation*}
\lambda = \frac{2z\mathsf{J}(g_J-1)^{2}}{ng_J^{2}\mu_{\mathrm{B}}^{2}}.\tag{5.18}
\end{equation*}
（对轨道淬灭的过渡金属离子特例 $g_J=2$，式 5.18 退化为式 5.16。）居里温度可写为
\begin{equation*}
T_{\mathrm{C}} = \frac{2z(g_J-1)^{2}\mathsf{J}}{3k_{\mathrm{B}}}J(J+1).\tag{5.19}
\end{equation*}
因此预期临界温度正比于 de Gennes 因子 $(g_J-1)^{2}J(J+1)$（其值亦列于表 5.2）。
de Gennes 因子最大的 Gd 是铁磁体，但稀土金属表现出各种不同的基态——包括反铁磁与
螺旋磁性，见下面几节。

\section{反铁磁性}

若交换相互作用为负（$\mathsf{J}<0$），分子场的取向使最近邻磁矩反平行排列有利——
这就是反铁磁性（antiferromagnetism）。它常出现在可以看作两套相互贯穿的子晶格的
体系中（见图 5.6）：一套子晶格上磁矩向上，另一套向下。图 5.6 中每个磁矩的最近邻
全部在另一套子晶格上。于是我们起初假定：一套子晶格上的分子场正比于另一套子晶格的
磁化强度，且无外加磁场。

\begin{figure}
\centering
\includegraphics[width=0.45\textwidth]{fig5_06.pdf}
\caption{反铁磁体可以分解为两套相互贯穿的子晶格。}
\end{figure}

\subsection{反铁磁体的外斯模型}

把“向上”子晶格标记为 +、“向下”子晶格标记为 $-$，则每套子晶格上的分子场为
\begin{equation*}
\begin{aligned}
B_{+} &= -|\lambda|M_{-}\\
B_{-} &= -|\lambda|M_{+},
\end{aligned}\tag{5.20}
\end{equation*}
$\lambda$ 是现在取负值的分子场常数。每套子晶格上
\begin{equation*}
M_{\pm} = M_{\mathrm{s}}\mathrm{B}_J\left(-\frac{g_J\mu_{\mathrm{B}}J|\lambda|M_{\mp}}{k_{\mathrm{B}}T}\right).\tag{5.21}
\end{equation*}
两套子晶格除磁矩方向外一切等价，故
\begin{equation*}
|M_{+}| = |M_{-}| \equiv M,\tag{5.22}
\end{equation*}
从而
\begin{equation*}
M = M_{\mathrm{s}}\mathrm{B}_J\left(\frac{g_J\mu_{\mathrm{B}}J|\lambda|M}{k_{\mathrm{B}}T}\right).\tag{5.23}
\end{equation*}
这与铁磁体的对应方程（式 5.6、5.7）几乎相同，每套子晶格上的分子场将严格遵循图 5.3
的形式，并在转变温度以上消失。该转变温度即奈尔温度（Néel temperature）$T_{\mathrm{N}}$：
\begin{equation*}
T_{\mathrm{N}} = \frac{g_J\mu_{\mathrm{B}}(J+1)|\lambda|M_{\mathrm{s}}}{3k_{\mathrm{B}}} = \frac{n|\lambda|\mu_{\mathrm{eff}}^{2}}{3k_{\mathrm{B}}}.\tag{5.24}
\end{equation*}
虽然每套子晶格的磁化强度都遵循图 5.3 的形式，但两者方向相反，反铁磁体的净磁化强度
$M_{+}+M_{-}$ 为零。可以定义错向磁化强度（staggered magnetization）为两套子晶格
磁化强度之差 $M_{+}-M_{-}$：它在 $T_{\mathrm{N}}$ 以下非零，因而可用作反铁磁体的
序参量（order parameter）。\fn{5}{序参量将在第 6.1 节定义。}

\subsection{磁化率}

$T>T_{\mathrm{N}}$ 时，小外场的效果可以与铁磁体同样计算——展开布里渊函数
$\mathrm{B}_J(y)=(J+1)y/3J+O(y^{3})$——得磁化率
\begin{equation*}
\chi = \lim_{B\to0}\frac{\mu_0M}{B} \propto \frac{1}{T+T_{\mathrm{N}}},\tag{5.25}
\end{equation*}
又是居里--外斯定律，只是 $-T_{\mathrm{C}}$ 换成了 $+T_{\mathrm{N}}$。

这给出解读顺磁态（即转变温度以上）磁化率数据的现成办法：把磁化率拟合为居里--外斯
形式
\begin{equation*}
\chi \propto \frac{1}{T-\theta},\tag{5.26}
\end{equation*}
$\theta$ 称为外斯温度。$\theta=0$：材料是顺磁体（见式 2.44）；$\theta>0$：材料是
铁磁体，预期 $\theta=T_{\mathrm{C}}$（见式 5.12）；$\theta<0$：材料是反铁磁体，
预期 $\theta=-T_{\mathrm{N}}$（见式 5.25）。这些可能性示于图 5.7。
\bio{Louis E. F. Néel\\（1904--2000）}

\begin{figure}
\centering
\includegraphics[width=0.32\textwidth]{fig5_07.pdf}
\caption{居里--外斯定律 $\chi\propto1/(T-\theta)$（$T>\theta$）：（a）三种情形——
$\theta=0$（顺磁体）、$\theta=\Theta>0$（铁磁体）与 $\theta=-\Theta<0$（反铁磁体）。
（b）$1/\chi$ 对 $T$ 作图为直线，温度轴截距给出 $\theta$。（c）$\chi T$ 对 $T$ 可为
常数（$\theta=0$）、随 $T$ 减小而增大（$\theta>0$）或随 $T$ 减小而减小
（$\theta<0$）。}
\end{figure}

实验测得的反铁磁体外斯温度常常离 $-T_{\mathrm{N}}$ 很远（一些常见反铁磁体的数据
见表 5.3）。这一偏差主要源于我们“一套子晶格上的分子场只依赖另一套子晶格磁化强度”
的假设。更现实的计算见习题 5.3。

\begin{table}
\centering
\caption{一些常见反铁磁体的性质。}
\begin{small}
\begin{tabular}{lccc}
\toprule
材料 & $T_{\mathrm{N}}$（K） & $\theta$（K） & $J$\\
\midrule
MnF$_2$ & 67 & $-80$ & $\frac{5}{2}$\\
MnO & 116 & $-510$ & $\frac{5}{2}$\\
CoO & 292 & $-330$ & $\frac{3}{2}$\\
FeO & 116 & $-610$ & 2\\
Cr$_2$O$_3$ & 307 & $-485$ & $\frac{3}{2}$\\
$\alpha$-Fe$_2$O$_3$ & 950 & $-2000$ & $\frac{5}{2}$\\
\bottomrule
\end{tabular}
\end{small}
\end{table}

在 $T_{\mathrm{N}}$ 以下对反铁磁体加磁场，比 $T_{\mathrm{C}}$ 以下的铁磁体复杂，
因为磁场的施加方向至关重要。磁矩沿场排列不再有能量优势：若两套子晶格磁化强度
等大反向，一套上省的能量会被另一套付出的能量抵消。

先考虑绝对零度（T=0），热扰动可以忽略：$|M_{+}|=|M_{-}|=M_{\mathrm{s}}$。若加
小磁场平行于一套子晶格的磁化方向（从而反平行于另一套），相当于给各子晶格的局域场
加上或减去一个小项；两套子晶格均已饱和，这不起作用，材料中诱导的净磁化强度为零，
故 $\chi_{\parallel}=0$。若小场垂直于磁化方向施加，则两套子晶格的磁化强度都略微
倾斜，沿外场产生一个磁化强度分量（见图 5.8），故 $\chi_{\perp}\neq0$。

\begin{figure}
\centering
\includegraphics[width=0.45\textwidth]{fig5_08.pdf}
\caption{$\chi_{\perp}$ 的起源。垂直于子晶格磁化方向加小磁场 B，两套子晶格的磁化
强度都略微倾斜，沿外场产生一个磁化强度分量。}
\end{figure}

温度升高但仍低于 $T_{\mathrm{N}}$ 时，热涨落减小每套子晶格上的分子场。这对平行
情形影响巨大——场使一套子晶格磁化增强、另一套减弱；而垂直情形下 $M_{+}$ 与
$M_{-}$ 等量减小，受小场的影响也对称，升温几乎没有影响。于是 $\chi_{\perp}$ 不
依赖温度，而 $\chi_{\parallel}$ 从 0 上升到 $T\to T_{\mathrm{N}}$ 时的
$\chi_{\perp}$（见图 5.9）。

\begin{figure}
\centering
\includegraphics[width=0.45\textwidth]{fig5_09.pdf}
\caption{温度对 $\chi_{\parallel}$ 与 $\chi_{\perp}$ 的影响。}
\end{figure}

\subsection{强磁场的效果}

先考虑 T=0 的反铁磁体在强磁场下的行为，以避开热涨落的干扰。场足够大时终将压倒
任何内部分子场，迫使所有磁矩彼此平行。但随场增大，虽然最终结局明确，通往终点的
路径却强烈依赖于外场相对子晶格磁化强度初始方向的取向。

若外场垂直于子晶格磁化，则随场增大，磁矩只是越来越弯（$\phi$ 逐渐变小，见
图 5.8），直到与外场排齐。

若外场平行于子晶格磁化，情形更有趣。小场下磁矩不转动、保持原位（图 5.10(a)）；
但达到临界场时，系统突然跳入另一种组态（图 5.10(b)），称为自旋翻转跳跃转变
（spin-flop transition）。场继续增大时角 $\theta$ 逐渐变小，最终磁矩与外场排齐。

\begin{figure}
\centering
\includegraphics[width=0.42\textwidth]{fig5_10.pdf}
\caption{磁场平行于子晶格磁化强度施加。（a）小场下无事发生，系统保持在反铁磁相。
（b）超过临界场后系统经自旋翻转跳跃转变进入自旋翻转跳跃相。}
\end{figure}

这些效应可以定量计算。设 $M_{+}$ 与场的夹角为 $\theta$（逆时针量），$M_{-}$ 与场
的夹角为 $\phi$（顺时针量），磁场沿晶体学 $z$ 轴施加。反铁磁相对应 $\theta=0$、
$\phi=\pi$（图 5.10(a)），自旋翻转跳跃相对应 $\theta=\phi$。需要判断哪个相能量
更低。

设总能量 E 是各子晶格塞曼能量与一个代表交换耦合（依赖两套子晶格磁矩相对取向）的
项之和：
\begin{equation*}
E = -MB\cos\theta - MB\cos\phi + AM^{2}\cos(\theta+\phi),\tag{5.27}
\end{equation*}
A 是与交换耦合有关的常数。为模拟磁各向异性，须加上形如
\begin{equation*}
-\frac{1}{2}\Delta(\cos^{2}\theta+\cos^{2}\phi)\tag{5.28}
\end{equation*}
的项（$\Delta$ 为小常数）——它反映磁化强度确实偏好沿某晶体学轴（此处 $z$ 轴）
排列：$\theta$、$\phi$ 偏好取 0 或 $\pi$ 而非中间值。反铁磁情形（$\theta=0$、
$\phi=\pi$）：$E=-AM^{2}-\Delta$，与场无关。自旋翻转跳跃情形（$\phi=\theta$）：
\begin{equation*}
E = -2MB\cos\theta + AM^{2}\cos2\theta - \Delta\cos^{2}\theta.\tag{5.29}
\end{equation*}
条件 $\partial E/\partial\theta=0$ 表明（忽略各向异性项）在 $\theta=\cos^{-1}[B/2AM]$
处有能量极小。代回 E 并作图得图 5.11：低于临界场 $B_{\text{spin-flop}}$ 时反铁磁相
能量最低；临界场处系统从一态切换到另一态，发生自旋翻转跳跃转变；高于此场时自旋
翻转跳跃相能量最低。

\begin{figure}
\centering
\includegraphics[width=0.45\textwidth]{fig5_11.pdf}
\caption{反铁磁相与自旋翻转跳跃相的能量随 B 的变化。}
\end{figure}

\begin{figure}
\centering
\includegraphics[width=0.45\textwidth]{fig5_12.pdf}
\caption{（a）对反铁磁体施加平行磁场时的磁化强度：起初无事发生，随后在 $B_1$ 处
发生自旋翻转跳跃转变进入自旋翻转跳跃相；磁场继续转动磁矩，直至 $B_2$ 处饱和。
（b）若自旋强烈偏好沿平行方向排列，则不发生自旋翻转跳跃，而是在 $B_3$ 处发生
自旋翻转（spin-flip）转变。两图均为绝对零度的预期曲线；有限温度会把尖角磨圆。
这也称为变磁转变（metamagnetic transition）。}
\end{figure}

大平行磁场下反铁磁体的磁化强度示于图 5.12(a)：自旋翻转跳跃转变之前毫无变化，其上
磁化强度稳步增大直至饱和。若各向异性效应很强（$\Delta$ 大），会出现另一种效应：
外场沿 z 时不再发生自旋翻转跳跃，而是发生自旋翻转（spin-flip）转变——B 达到临界
值时一套子晶格的磁化强度突然反转，系统一步进入铁磁态（图 5.12(b)）。

\subsection{反铁磁有序的类型}

反铁磁性的另一个复杂之处：把等量的上下自旋排在一个晶格上有很多种方式，且不同
排列还取决于晶格类型。图 5.13、5.14 给出若干可能的排列。

立方钙钛矿中磁性原子排在简单立方晶格上，G 型有序（图 5.13(d)）非常常见——经氧
原子的超交换相互作用迫使所有最近邻磁性原子反铁磁排列。这只对 G 型有序成立，
例如 LaFeO$_3$ 与 LaCrO$_3$。LaMnO$_3$ 也是立方钙钛矿，却表现 A 型有序
（图 5.13(a)）：铁磁排列的 (100) 面交替取向。原因是 $\mathrm{Mn^{3+}}$ 离子的
扬--特勒畸变使 (100) 面内 Mn--O 键长短交替，相邻 $\mathrm{Mn^{3+}}$ 离子的轨道
取向不同，超交换导致一个原子的占据轨道与邻居的未占据轨道之间的相互作用：面内
相互作用为铁磁，面间则因常规超交换而为反铁磁。

\begin{figure}
\centering
\includegraphics[width=0.24\textwidth]{fig5_13.pdf}
\caption{简单立方晶格上可以出现的四种反铁磁有序（A、C、E、G 型）。两种自旋态以
+、$-$ 标记。}
\end{figure}

\begin{figure}
\centering
\includegraphics[width=0.28\textwidth]{fig5_14.pdf}
\caption{体心立方晶格上可以出现的三种反铁磁有序。}
\end{figure}

\section{亚铁磁性}

以上对反铁磁性的处理假定两套子晶格等价。若因晶体学上的原因二者不等价呢？此时
两套子晶格的磁化强度未必等大反向，因而不相抵消——材料将有净磁化强度。这一现象
称为亚铁磁性（ferrimagnetism）。由于两套子晶格上的分子场不同，子晶格的自发磁化
强度对温度的依赖可以大相径庭，净磁化强度本身因此可能有复杂的温度依赖。有时低温
下一套子晶格主导磁化强度，较高温度时另一套主导；此时净磁化强度可在某个温度降至
零并变号，该温度称为补偿温度（compensation temperature）。因此亚铁磁体的磁化率
不遵循居里--外斯定律。

铁氧体（ferrites）是一族亚铁磁体：化学式 $MO\cdot\mathrm{Fe_2O_3}$ 的化合物，
M 是二价阳离子（如 $\mathrm{Zn^{2+}}$、$\mathrm{Co^{2+}}$、$\mathrm{Fe^{2+}}$、
$\mathrm{Ni^{2+}}$、$\mathrm{Cu^{2+}}$ 或 $\mathrm{Mn^{2+}}$）。晶体结构为尖晶石
结构（spinel），含两类格位：四面体位（四个氧邻居，称 A 位）与八面体位（六个氧
邻居，称 B 位）；B 位数目是 A 位两倍。两套子晶格不等价——存在两类晶体学格位，
且容纳两种不同离子。正尖晶石（normal spinel）中 $M^{2+}$ 阳离子坐 A 位，
$\mathrm{Fe^{3+}}$（${}^{6}S_{5/2}$，磁矩 $5\mu_{\mathrm{B}}$）阳离子坐 B 位。
反尖晶石（inverse spinel）中 $M^{2+}$ 阳离子坐一半 B 位，$\mathrm{Fe^{3+}}$ 阳离子
占据另一半 B 位与全部 A 位；反尖晶石中 A、B 位上 $\mathrm{Fe^{3+}}$ 阳离子的磁矩
反平行，样品总磁矩仅来自 $M^{2+}$ 离子。

\begin{table}
\centering
\caption{一些常见亚铁磁体的性质。}
\begin{small}
\begin{tabular}{lccc}
\toprule
材料 & $T_{\mathrm{C}}$（K） & 磁矩（$\mu_{\mathrm{B}}$/化学式单元） & 补偿温度（K）\\
\midrule
Fe$_3$O$_4$ & 858 & 4.1 & --\\
CoFe$_2$O$_4$ & 793 & 3.7 & --\\
NiFe$_2$O$_4$ & 858 & 2.3 & --\\
CuFe$_2$O$_4$ & 728 & 1.3 & --\\
Y$_3$Fe$_5$O$_12$ & 560 & 5.0 & --\\
Gd$_3$Fe$_5$O$_12$ & 564 & 16.0 & 290\\
Dy$_3$Fe$_5$O$_12$ & 563 & 18.2 & 220\\
Ho$_3$Fe$_5$O$_12$ & 567 & 15.2 & 137\\
\bottomrule
\end{tabular}
\end{small}
\end{table}

M=Fe 的情形即 $\mathrm{Fe_3O_4}$（半导体，与其他绝缘体铁氧体不同），已在 4.2.5 节
讨论过。

另一族亚铁磁体是石榴石（garnets），化学式 $R_3\mathrm{Fe_5O_{12}}$，R 是三价稀土
原子。晶体结构为立方，但单胞相当复杂：三个 $\mathrm{Fe^{3+}}$ 离子在四面体位，
两个在八面体位，$R^{3+}$ 离子在十二面体对称性的格位上。钇铁石榴石（YIG，
$\mathrm{Y_3Fe_5O_{12}}$）中 $Y^{3+}$ 无磁矩（$4d^{0}$），四面体位 $\mathrm{Fe^{3+}}$
的磁矩与八面体位的反平行，净磁矩为 $5\mu_{\mathrm{B}}$。

钡铁氧体（$\mathrm{BaFe_{12}O_{19}}=\mathrm{BaO}\cdot6\mathrm{Fe_2O_3}$）具有
六角结构：八个 $\mathrm{Fe^{3+}}$ 离子与其他四个反平行，净磁矩相当于四个
$\mathrm{Fe^{3+}}$ 离子，即 $20\mu_{\mathrm{B}}$。粉体形式下它用于磁记录，因为
矫顽力高（见 6.7.9 节）。一些常见亚铁磁体的性质列于表 5.4。

大多数亚铁磁体是电绝缘体，这一事实造就了它们的许多实际应用。铁磁体往往是金属，
不适合涉及振荡磁场的应用：快速变化的磁场感应电压并在导体中驱动电流（涡流，
eddy currents），电流在金属中造成电阻性发热（涡流损耗）。因此，当需要自发磁化
材料工作在高频时，常可用许多亚铁磁体——绝缘体中感应电压无法驱动显著的涡流。
固体铁氧体磁芯用于许多高频应用：需要高磁导率、低能量损耗的天线与变压器，以及
微波器件。此外，许多亚铁磁体本就是氧化物，比金属铁磁体更耐腐蚀。

\section{螺旋有序}

许多稀土金属的晶体结构使原子分层排列。先考虑与镝相关的情形：层内原子磁矩铁磁
排列，层间相互作用由最近邻交换常数 $\mathsf{J}_1$ 与次近邻交换常数 $\mathsf{J}_2$
描述。若相邻两基面（即层所在平面）中磁矩间夹角为 $\theta$（见图 5.15(a)），系统
能量可写为
\begin{equation*}
E = -2NS^{2}(\mathsf{J}_1\cos\theta + \mathsf{J}_2\cos2\theta),\tag{5.30}
\end{equation*}
N 是每层原子数。能量在 $\partial E/\partial\theta=0$ 时取极小：
\begin{equation*}
(\mathsf{J}_1 + 4\mathsf{J}_2\cos\theta)\sin\theta = 0.\tag{5.31}
\end{equation*}
解为 $\sin\theta=0$（即 $\theta=0$ 或 $\theta=\pi$，铁磁或反铁磁），或
\begin{equation*}
\cos\theta = -\frac{\mathsf{J}_1}{4\mathsf{J}_2}.\tag{5.32}
\end{equation*}
最后这个解对应螺旋有序（helical order，又称螺旋磁性 helimagnetism），当
$\mathsf{J}_2<0$ 且 $|\mathsf{J}_1|<4|\mathsf{J}_2|$ 时优于铁磁或反铁磁
（见图 5.15(b)）。螺旋的螺距一般与晶格常数不可公度，故晶体中没有两层恰好具有
相同的自旋方向。

\begin{figure}
\centering
\includegraphics[width=0.4\textwidth]{fig5_15.pdf}
\caption{（a）螺旋磁有序。（b）由最近邻交换常数 $\mathsf{J}_1$ 与次近邻交换常数
$\mathsf{J}_2$ 耦合的层模型的相图。}
\end{figure}

螺旋结构出现在许多磁性体系中，最著名的是稀土金属。许多稀土金属具有六角密堆晶体
结构；螺旋轴垂直于六角密堆面，沿通常定义的 c 轴。Tb、Dy、Ho 中自旋旋转的平面是
六角密堆面；而 Er、Tm 中易轴为 c 轴，故在某温度范围内自旋的 c 分量被正弦调制。
还可能出现螺形结构（图 5.1(d)）：自旋有沿易轴（垂直于旋转面）的分量，同时面内
分量绕螺旋进动。

稀土金属中的交换相互作用是由传导电子传递的间接 RKKY 相互作用：比超交换长程得多，
且随距离变号。上述只含最近邻与次近邻相互作用的模型因此过分简化——计算波矢依赖
的交换相互作用 $\mathsf{J}(\mathbf{q})$ 需要每个稀土金属费米面的细节；若它在某
波矢 q 取极大，则可诱发波矢为 q 的螺旋磁性。另一个因素是稀土金属中晶体场的大
影响：晶体场分裂电子态，通常不仅基态、某些激发态也有热布居——使理解这些体系
所需的分析更加复杂。

\section{自旋玻璃}

迄今考虑的材料每个单胞里都有一个或多个磁矩。若从一个非磁晶格出发，稀疏地、随机
地散布磁性原子，会怎样？直觉或许会认为：结局是彻底随机的东西，不大可能表现出从高温
无序态到低温有序态的相变。这只对了一部分：这类体系虽然内在随机，却在某个特定
温度表现出近似于相变的行为，进入一个虽非有序、却与高温无序态截然不同的态。自旋
玻璃（spin glass）可定义为：具有混合相互作用的随机磁性系统，其特征是在明确的环境
温度 $T_{\mathrm{f}}$（冻结温度）以下自旋发生随机然而协同的冻结，冻结态是亚稳的，
且没有通常的磁长程序。

\begin{example}{5.1}
研究得很好的自旋玻璃例子是 CuMn 合金，Mn 浓度为几个原子百分比。Mn 离子稀疏，
其磁矩经 Cu 中传导电子传递的 RKKY 相互作用（见 4.2.4 节与 7.7.3 节）彼此耦合。
RKKY 相互作用随符号振荡，相互作用可铁磁可反铁磁。净结果是一个内置大量阻挫
（frustration）、没有明确基态而有大量可能基态的系统。高温下 Mn 磁矩热涨落；温度
降低时磁矩变慢、局域关联的团簇逐渐建立。在 $T_{\mathrm{f}}$ 处磁矩卡进这些简并
基态之一，冻结。CuMn 中位置随机性给出自旋间距离分布，从而提供阻挫；自旋玻璃也
见于具有键随机性（bond-randomness）的体系——最近邻相互作用在 $+\mathsf{J}$ 与
$-\mathsf{J}$ 之间变化。自旋玻璃与阻挫将在 8.2 节更详细讨论。
\end{example}

\section{核有序}

迄今只考虑了电子磁矩。有些核具有磁矩——它们能有序吗？答案是能，但效应极小，
只在低温且材料没有电子磁矩时才重要。核自旋深居原子内部，彼此耦合不强，其间不
存在直接交换相互作用。核磁矩比电子磁矩小约三个量级，而偶极相互作用正比于磁矩
的平方，故任何偶极相互作用都比电子情形低六个量级。只要电子磁矩存在，任何核效应
都完全被淹没。核磁矩之小意味着核自旋体系只能在低于 $\sim1\ \mu$K 的温度显示磁
有序。金属样品中核磁矩之间也存在 RKKY 耦合的可能——交换由传导电子传递——但其
大小与核偶极相互作用同量级。

尽管如此，在若干没有电子磁矩的材料上已做出惊人的实验，核自旋有序真的被观察到：
铜在 58 nK 发生一级相变进入反铁磁有序的核基态；银中核自旋的奈尔温度为 560 pK。
这些极低温度如何获得？实验依靠一个事实：同一块样品中，晶格与核可以温度不同。
用绝热退磁（见 2.6 节）冷却核，核在由自旋--自旋弛豫时间 $T_2$（典型几毫秒）
表征的时间内达到彼此的热平衡；而核自旋要以由 $T_1$（可长达数小时）表征的时间才
弛豫回晶格温度。

还有一个巧妙的窍门：把核置于磁场中，核能级分裂并按玻尔兹曼概率布居。然后把磁场
突然反向（时间短于 $T_2$）：核自旋来不及相对能级移动，于是得到布居反转；而且
各能级之间恰具有对应负温度的正确玻尔兹曼分布。负温度态是一个有着独特磁相图的
独特热力学态。银的实验表明：低温负温度下银核在 $-1.9$ nK 发生铁磁有序。

\section{磁有序的测量}

前面几节描述了各种磁有序。但如何测量？本节讨论测量磁有序的各种技术。

\subsection{磁化强度与磁化率}

最直截了当的实验技术是常规磁化强度测量——本质上测样品的净磁矩 $\boldsymbol{\mu}$
（除以样品体积即得磁化强度）。

抽拉式磁强计（extraction magnetometer）：样品置于线圈中心，然后抽出到很远处，在线圈中
感应电压 V；样品产生的磁通 $\Phi$ 等于 $\int V\,\mathrm{d}t$。通常磁强计由两个
反绕线圈构成，使外场变化引起的电压相互抵消，只留下样品的信号。振动样品磁强计
（vibrating sample magnetometer，通称 VSM）让样品正弦上下振动，样品磁矩的运动在
固定的拾取线圈中感应电信号；信号正比于磁矩以及振动幅度与频率。交变梯度磁强计中，
样品贴在压电条上，反绕线圈提供交变场，使样品处出现交变的场梯度，样品受力
$(\boldsymbol{\mu}\cdot\nabla)\mathbf{B}$；压电条以驱动频率同频偏折，偏折量可由
压电信号测得，从而推知 $\boldsymbol{\mu}$。力矩磁强计中样品悬在扭丝上，加磁场 B
产生力矩 $\boldsymbol{\mu}\times\mathbf{B}$。最灵敏的技术之一用 SQUID（超导量子
干涉器件）：一个含“弱连接”（即约瑟夫森结）的超导环，像一个极灵敏的量子干涉仪；
让样品穿过环，感应的持续电流正比于样品磁化强度。这些技术当然只在样品磁化强度
非零时有用；反铁磁体的磁化强度在外场不大时为零。

另一策略是测磁化率——它很好地指示磁有序的类型，可拟合成居里--外斯定律。
磁化率有多种测法。它等于小外场下样品中诱导的磁化强度，故上述测磁化强度的技术
皆可用于测磁化率。此外，主要用于提取小磁化率的方法是古埃法（Gouy method）：
长圆柱样品悬于磁体极靴之间，一端在极靴间、另一端远在磁体之外零场区。通电后样品
的表观重量改变。相关的法拉第法可用更小的样品，但此时力由适当形状极靴产生的磁场
梯度提供。

\subsection{中子散射}

事实证明，中子在研究凝聚态磁性中极为有用；在研究磁结构与动力学方面，中子散射比其他
实验技术有显著优势：它给出样品中磁矩排列的最直接信息。

中子是自旋 $\frac{1}{2}$ 粒子（见表 2.3），具有非零磁矩。核反应堆燃料元件内的
裂变反应产生大量中子。反应堆出来的中子束的能谱由慢化剂温度 T 决定——通常为室温，
加热或冷却慢化剂可改变 T。因此反应堆中子具有速度分布
\begin{equation*}
n(v) \propto v^{3}\exp\left(-\frac{\frac{1}{2}m_{\mathrm{n}}v^{2}}{k_{\mathrm{B}}T}\right),\tag{5.33}
\end{equation*}
$n(v)\,\mathrm{d}v$ 是每秒通过单位面积的速度在 $v$ 到 $v+\mathrm{d}v$ 的中子数，
$m_{\mathrm{n}}$ 是中子质量。分布含玻尔兹曼因子与 $v^{3}$ 项——后者是相空间因子
$v^{2}$（同于分子运动论的麦克斯韦速度分布）与泻流过孔的因子 v 之积。此函数绘于
图 5.16(a)，其极大在
\begin{equation*}
v = \sqrt{\frac{3k_{\mathrm{B}}T}{m_{\mathrm{n}}}}\tag{5.34}
\end{equation*}
对应条件 $\frac{1}{2}m_{\mathrm{n}}v^{2}=\frac{3}{2}k_{\mathrm{B}}T$。速度为 v 的
中子的德布罗意波长 $\lambda$ 为
\begin{equation*}
\lambda = \frac{h}{m_{\mathrm{n}}v}.\tag{5.35}
\end{equation*}
分布函数按波长绘于图 5.16(b)。这表明热中子的波长与原子间距相近，从而可以进行
衍射测量。波长可调意味着衍射实验的尺度范围从直接探测小原子的波函数直到大分子的
低分辨研究。

\begin{figure}
\centering
\includegraphics[width=0.42\textwidth]{fig5_16.pdf}
\caption{（a）30 K 与 300 K 时中子分布函数随速度 v 的变化。（b）同样的分布函数按
德布罗意波长 $\lambda=h/m_{\mathrm{n}}v$ 绘出。曲线已归一化为等面积。}
\end{figure}

中子经强核力与核散射，也经电磁相互作用与晶体内磁场的变化散射。后者使中子能探测
材料磁性——中子与电子的磁矩耦合。

先考虑核散射。引起中子散射的核力程极短（$10^{-15}$--$10^{-14}$ m），远小于典型
热中子波长。因此入射中子波 $\psi_i=e^{i\mathbf{k}\cdot\mathbf{r}}$ 产生球对称的
弹性散射波 $\psi_f=-(b/r)e^{ikr}$。量 b 称为散射长度（scattering length）；表达式
中的负号保证排斥性核势对应 $b>0$。散射长度依赖具体的核以及核--中子系统的自旋态
（对自旋 I 的核，可为 $I+\frac{1}{2}$ 或 $I-\frac{1}{2}$）。

弹性中子散射实验中，弹性散射的中子在散射矢量等于倒格矢时产生强布拉格反射
（图 5.17(a, b)）。研究这些反射的实验方法包括：（i）转动晶体并测量单色束的散射，
寻找强反射；（ii）使用一系列入射波长（劳厄法，见图 5.17(c)）；（iii）测量单色
中子在粉末样品上的衍射（见图 5.17(d)）。反应堆源上一种可能的实验安排示意于
图 5.17(e)。中子也可由散裂源产生：同步加速器的高能质子轰击重金属靶（如
${}^{238}$U）产生中子；同步加速器质子束是脉冲的，中子束因此也是脉冲的。利用脉冲
中子的能量展宽可建造飞行时间谱仪：测量脉冲开始后散射中子随时间的变化，一次脉冲
即可提取一系列入射波矢的散射（见图 5.17(f)）。

中子的核散射来自与核（而非电子云）的相互作用，因此原子的散射能力与其原子序数
关系不大——这与 X 射线散射和电子散射截然相反。这带来若干优点：更容易在重原子
存在下感知轻原子（如氢）；周期表中相邻元素的散射截面一般差别显著、可以区分；
核散射的依赖性使同一元素的不同同位素可以有很不一样的中子散射长度，因此可以采用
同位素替换。中子是高度穿透的探针，能（同样是无损地）研究材料内部，而不像 X 射线
散射、电子显微术或光学方法那样仅研究表面层。

中子不带电——这正是电子云对它不可见的原因。但若原子有净磁矩，中子的磁矩可直接
与其耦合。这是中子重要的第二类散射机制：磁散射。结果是自旋向上与向下的电子磁矩
在中子“眼”里不同：样品变为磁性时，中子衍射图中可以出现新的峰。例子见图 5.18、
5.19——反铁磁体 MnO 的磁结构。$T_{\mathrm{N}}=116$ K 以上，中子衍射峰来自锰离子的
面心立方（fcc）晶格——所有锰离子等价（见图 5.19）。fcc 晶格使衍射图出现系统性
消光。$T_{\mathrm{N}}$ 以下，磁结构如图 5.18：每个 (111) 面内自旋全部平行，但相邻
两个 (111) 面的磁矩反平行。磁性单胞因而是化学单胞的两倍大，故 $T_{\mathrm{N}}$
以下中子衍射图中突然多出几个峰。磁布拉格峰的幅度可用来量度磁有序的强度，从而
可以跟踪磁有序随温度的变化。

\begin{figure}
\centering
\includegraphics[width=0.45\textwidth]{fig5_17.pdf}
\caption{（a）波矢 $\mathbf{k}$ 的入射中子被晶面以散射角 $2\theta$ 散射到
$\mathbf{k}'$；弹性散射意味着 $|\mathbf{k}|=|\mathbf{k}'|$；入射中子波长
$\lambda=2\pi/|\mathbf{k}|$。（b）当散射矢量 $\mathbf{Q}=\mathbf{k}-\mathbf{k}'$
等于倒格矢时，晶面散射的中子发生相长干涉。（c）劳厄法：入射方向固定而波长有一
范围的径束，当 $\mathbf{Q}$ 为倒格矢时给出布拉格反射。（d）粉末衍射法：单色束
入射粉末样品，散射中子位于半角 $2\theta$ 的圆锥（德拜--谢乐锥）上。（e）反应堆源
的中子衍射：单色器与旋转屏蔽选择入射中子波长，样品散射的中子由多探测器测量。
（f）散裂源：靶上散裂产生中子脉冲；探测器信号按时间记录，较快（波长较短）的
中子较早到达。}
\end{figure}

\begin{figure}
\centering
\includegraphics[width=0.42\textwidth]{fig5_18.pdf}
\caption{MnO 的磁结构。Mn 离子位于面心立方晶格上；$\mathrm{O^{2-}}$ 离子未画出
（见图 4.2）。Mn 离子按自旋态以黑白两色标记。}
\end{figure}

\begin{figure}
\centering
\includegraphics[width=0.5\textwidth]{fig5_19.pdf}
\caption{$T_{\mathrm{N}}$ 上、下 MnO 的中子衍射图。取自 C.~G.~Shull,
W.~A.~Strauser 与 E.~O.~Wollan, Phys.~Rev., \textbf{83}, 333 (1951)。}
\end{figure}

磁矩散射的幅度依赖于磁矩排列的方向，因此中子衍射可用于确定磁有序晶体中原子磁矩
的排列。

周期性体系中，散射矢量 Q 等于倒格矢 G 时出现尖锐布拉格反射。螺旋磁体中，螺旋
带来额外的周期性，可用螺旋波矢 q 描述：螺旋轴沿 q，螺距等于 $2\pi/|q|$。布拉格
反射于是也出现在 $G\pm q$ 处。由于螺旋螺距大于晶格间距，$|q|<|G|$：磁转变温度
以上观察到的布拉格峰在转变以下仍被观察到，且每个布拉格峰两侧伴随较小的卫星峰。

中子散射是确定样品磁态细节的最直接方法。但中子昂贵，且中子与物质相互作用较弱，
需要大样品（这一缺点对 6.6.4 节描述的非弹性实验最为严重）。

\subsection{其他技术}

内场可用第三章描述的技术测量。NMR、穆斯堡尔与 $\mu$SR 实验都能测铁磁体磁化强度
的温度依赖。它们是局域探针——测的不是平均磁化强度而是特定晶体学格位处的磁化
强度——因此可用于测量反铁磁体或亚铁磁体中各子晶格的磁化强度，还能给出样品整体
未磁化时磁畴内部自发磁化强度的信息。

另一种技术是 X 射线散射。磁性 X 射线散射通常是相当弱的效应（衍射实验中磁散射
名义上比非磁散射小五个量级），因此 X 射线以往主要用于结构而非磁性研究。这一局面
已因 X 射线共振交换散射（XRES）的发展而改变——当 X 射线能量扫过吸收边时磁散射
被增强。另一个重要因素是高亮度 X 射线同步辐射源的发展与优质偏振光子束的提供。
这些因素抵消了效应的内在弱点，使 X 射线得以用于磁性研究。X 射线的优点：可研究
小样品（束径可以很小）；元素分辨（把能量调到样品中特定原子的吸收边，测到的信号
就是该元素的标志）；分辨率可优于中子；还能分别测量磁化密度中自旋与轨道的贡献。

二色性（dichroism）是光的吸收对偏振的依赖（由取向长链分子构成的偏振片是光频线
二色性的著名例子）。磁二色性以类似方式工作，其效应由未成对电子造成的原子周围
磁化分布的不对称驱动。这本质上是轨道现象——光子把螺旋度 $\pm1$ 传递给吸收体——
但自旋--轨道相互作用可使二色性依赖自旋角动量。磁 X 射线圆二色（MXCD）技术测量
右旋与左旋圆偏振 X 射线衰减之差，结果可分别推知轨道与自旋磁矩。

\section*{延伸阅读}

\begin{itemize}[itemsep=0.2em]
\item 核磁有序实验的综述见 A.~S.~Oja 与 O.~V.~Lounsmaa, Reviews of Modern Physics
\textbf{67}, 1 (1997)。
\item 中子散射的描述见 M.~Dove,《Structure and dynamics》, OUP（即将出版）；
G.~L.~Squires,《Introduction to the theory of thermal neutron scattering》,
CUP 1978；W.~Marshall 与 S.~W.~Lovesey,《Theory of thermal neutron scattering》,
OUP 1971。
\item 更进阶的处理见 S.~W.~Lovesey,《Theory of neutron scattering from condensed
matter》, OUP 1984。
\item S.~W.~Lovesey 与 S.~P.~Collins,《X-ray scattering and absorption by magnetic
materials》, OUP 1996，提供了磁 X 射线散射的有用信息。
\end{itemize}

\section*{习题}

\exer{5.1} 估算铁中分子场 $B_{\mathrm{mf}}$ 的大小（以特斯拉计），并与磁化强度对
B 场的贡献 $\mu_0M$ 比较。进而说明为什么解释铁的铁磁性必须借助交换的概念。Fe 的
密度 7873 kg\,m$^{-3}$，相对原子质量 55.847，居里温度 $T_{\mathrm{C}}=1043$ K，
每个原子携带约 2.2 $\mu_{\mathrm{B}}$ 的磁矩。

\exer{5.2} 把外斯模型推广到自旋 $S>\frac{1}{2}$。证明恰低于 $T_{\mathrm{C}}$ 处
磁化强度为
\begin{equation*}
\frac{M}{M_{\mathrm{s}}} = \left(\frac{10(S+1)^{2}(T_{\mathrm{C}}-T)}{3[(S+1)^{2}+S^{2}]T_{\mathrm{C}}}\right)^{1/2}\tag{5.36}
\end{equation*}
且 $T_{\mathrm{C}}$ 处热容有大小为
\begin{equation*}
\Delta C = \frac{5}{2}nk_{\mathrm{B}}\frac{(2S+1)^{2}-1}{(2S+1)^{2}+1}.\tag{5.37}
\end{equation*}
的间断。你需要关系式
\begin{equation*}
\mathrm{B}_S(y) = \left[\frac{(2S+1)^{2}-1}{(2S)^{2}}\right]\frac{y}{3} - \left[\frac{(2S+1)^{4}-1}{(2S)^{4}}\right]\frac{y^{3}}{45} + O(y^{5}),\qquad y\ll1.\tag{5.38}
\end{equation*}

\exer{5.3} 反铁磁体每套子晶格上的分子场为
\begin{equation*}
\begin{aligned}
B_{+} &= -\Gamma M_{+} - |\lambda|M_{-}\\
B_{-} &= -\Gamma M_{-} - |\lambda|M_{+},
\end{aligned}\tag{5.39}
\end{equation*}
$\Gamma$ 是表示同一子晶格对分子场贡献的常数。推广 5.2 节的处理，证明每套子晶格上
磁化强度为
\begin{equation*}
M_{\pm} = M_{\mathrm{s}}\mathrm{B}_J\left(-\frac{g_J\mu_{\mathrm{B}}J(\Gamma M_{\pm}+|\lambda|M_{\mp})}{k_{\mathrm{B}}T}\right).\tag{5.40}
\end{equation*}
进而证明奈尔温度（每套子晶格自发磁化强度消失的温度）为
\begin{equation*}
T_{\mathrm{N}} = \frac{n(|\lambda|-\Gamma)\mu_{\mathrm{eff}}^{2}}{3k_{\mathrm{B}}}.\tag{5.41}
\end{equation*}
再证明磁化率为 $\chi\propto(T-\theta)^{-1}$，其中
\begin{equation*}
\theta = -\frac{n(|\lambda|+\Gamma)\mu_{\mathrm{eff}}^{2}}{3k_{\mathrm{B}}},\tag{5.42}
\end{equation*}
故仅当 $\Gamma=0$ 时 $|\theta|=T_{\mathrm{N}}$。

\exer{5.4} 对 5.4 节描述的螺旋磁体，能量由式 5.30 给出。证明铁磁、反铁磁与螺旋磁
情况的能量分别为 $E_{\mathrm{FM}}$、$E_{\mathrm{AFM}}$ 与 $E_{\mathrm{HM}}$：
\begin{equation*}
E_{\mathrm{FM}} = -2NS^{2}(\mathsf{J}_1+\mathsf{J}_2)\tag{5.43}
\end{equation*}
\begin{equation*}
E_{\mathrm{AFM}} = -2NS^{2}(-\mathsf{J}_1+\mathsf{J}_2)\tag{5.44}
\end{equation*}
\begin{equation*}
E_{\mathrm{HM}} = -2NS^{2}\left(\frac{-\mathsf{J}_1^{2}}{8\mathsf{J}_2} - \mathsf{J}_2\right)\tag{5.45}
\end{equation*}
故当 $\mathsf{J}_2<0$ 且 $|\mathsf{J}_1|<4|\mathsf{J}_2|$ 时螺旋磁性优于铁磁或
反铁磁。

\exer{5.5} 证明 $T\ll T_{\mathrm{C}}$ 时铁磁体的外斯模型给出磁化强度
\begin{equation*}
M = M_{\mathrm{s}}(1-2\,\mathrm{e}^{-2T_{\mathrm{C}}/T}).\tag{5.46}
\end{equation*}
注意这与实验确定的幂律行为不符——实验给出
\begin{equation*}
M = M_{\mathrm{s}}(1-AT^{3/2})\tag{5.47}
\end{equation*}
A 为常数。这一低温行为可以由自旋波效应解释（见 6.6 节）。

\exer{5.6} 亚铁磁体的磁化率 $\chi$ 可用平均场理论处理。考虑具有两套不等价子晶格的
亚铁磁体，子晶格 1、2 上的分子场 $B_1$、$B_2$ 为
\begin{equation*}
B_1 = \mu_0H - \lambda M_1\tag{5.48}
\end{equation*}
\begin{equation*}
B_2 = \mu_0H - \lambda M_2,\tag{5.49}
\end{equation*}
$M_1$、$M_2$ 是各子晶格的磁化强度，H 是外场。但两套子晶格有不同的居里常数
$C_1$、$C_2$（即不同的 $g_JJ$ 值）。证明磁化率为
\begin{equation*}
\chi = \frac{1}{T^{2}-\theta^{2}}\left[(C_1+C_2)T - \frac{2\lambda C_1C_2}{\mu_0}\right],\tag{5.50}
\end{equation*}
并求转变温度 $\theta$ 的值。证明当 $C_1=C_2=C$ 时磁化率退化为 $\theta=\lambda C/\mu_0$
的反铁磁体情形。

\exer{5.7} 证明海森伯模型的哈密顿量
\begin{equation*}
\hat{\mathcal{H}} = -\sum_{\langle ij\rangle}\mathsf{J}\,\mathbf{S}_i\cdot\mathbf{S}_j,\tag{5.51}
\end{equation*}
可改写为
\begin{equation*}
\hat{\mathcal{H}} = -\sum_{\langle ij\rangle}\mathsf{J}\left[S_i^{z}S_j^{z} + \frac{1}{2}(S_i^{+}S_j^{-}+S_i^{-}S_j^{+})\right].\tag{5.52}
\end{equation*}
进而证明：对海森伯铁磁体（$\mathsf{J}>0$），所有格点自旋都朝上（比如说）的态
$|\Phi\rangle$ 是哈密顿量的本征态，能量为
\begin{equation*}
E_0 = -NS^{2}\mathsf{J}.\tag{5.53}
\end{equation*}
考虑海森伯反铁磁体（$\mathsf{J}<0$），自旋位于两套子晶格上，每个自旋只与
另一套子晶格上的自旋相互作用。证明“想当然”的基态——每套子晶格内铁磁排列、两套
子晶格磁化方向相反——并非哈密顿量的本征态。这凸显了海森伯反铁磁体是一个复杂
而困难的问题。

\exer{5.8} $\mathrm{MnF_2}$ 具有四角晶体结构：Mn 离子位于四角单胞（a=b=0.5 nm，
c=0.3 nm）的顶角及体心位置。70 K 以下 Mn 离子自旋沿 c 轴反铁磁排列，体心离子的
自旋与顶角离子的相反。用入射波长 0.3 nm、散射角 0°--90° 测量 $\mathrm{MnF_2}$
粉末样品的中子散射。画出（a）100 K 与（b）10 K 时预期的结果。可忽略 F 离子的
散射。

\exer{5.9} 外斯模型可以找到精确解。无外场时问题归结为联立求解
\begin{equation*}
\tilde{m} = \frac{M}{M_{\mathrm{s}}} = \mathrm{B}_J(y)\tag{5.54}
\end{equation*}
与
\begin{equation*}
y = \frac{g_J\mu_{\mathrm{B}}J\lambda M}{k_{\mathrm{B}}T}\tag{5.55}
\end{equation*}
$\tilde{m}$ 是约化磁化强度，$M_{\mathrm{s}}=ng_J\mu_{\mathrm{B}}J$。定义 $x=e^{y}$
有助益。

（a）对 $J=\frac{1}{2}$、$g_J=2$，$\mathrm{B}_J(y)=\tanh y$。证明这意味着
\begin{equation*}
\tilde{m} = \frac{x-\frac{1}{2}}{x+\frac{1}{2}}\tag{5.56}
\end{equation*}
从而
\begin{equation*}
\log x = \frac{1}{2}\left[\log(1+\tilde{m})-\log(1-\tilde{m})\right].\tag{5.57}
\end{equation*}
写 $T_{\mathrm{C}}=n\mu_{\mathrm{B}}^{2}\lambda/k_{\mathrm{B}}$，证明
\begin{equation*}
\frac{T}{T_{\mathrm{C}}} = \frac{2\tilde{m}}{\log(1+\tilde{m})-\log(1-\tilde{m})}.\tag{5.58}
\end{equation*}

（b）对 $J=1$，证明
\begin{equation*}
\mathrm{B}_1(y) = \frac{2\sinh y}{2\coth y+1}\tag{5.59}
\end{equation*}
从而
\begin{equation*}
\tilde{m} = \frac{x-\frac{1}{x}}{x+1+\frac{1}{x}}\tag{5.60}
\end{equation*}
其中
\begin{equation*}
x = \exp(2\lambda M\mu_{\mathrm{B}}/k_{\mathrm{B}}T).\tag{5.61}
\end{equation*}
证明这意味着
\begin{equation*}
x = \frac{\tilde{m}\pm\sqrt{4-3\tilde{m}^{2}}}{2(1-\tilde{m})},\tag{5.62}
\end{equation*}
进而（仔细想想 $\pm$ 号）意味着
\begin{equation*}
\frac{T}{T_{\mathrm{C}}} = \frac{\frac{3}{2}\tilde{m}}{\log(\tilde{m}+\sqrt{4-3\tilde{m}^{2}})-\log(2(1-\tilde{m}))},\tag{5.63}
\end{equation*}
其中 $T_{\mathrm{C}}=8n\mu_{\mathrm{B}}^{2}\lambda/3k_{\mathrm{B}}$。本题所用方法
详述于 M.~A.~B.~Whittaker, American Journal of Physics \textbf{57}, 45 (1989)。
